\scanpage{438}
\mspage{350}
\uline{Lemme} 4.13.30. Soit $I$ une catégorie grde devant $\pi$, et $(J_{\unsure{r|i}})_{r\in A}$ une famille de sous-catégories pleines telles que $\mathrm{Card}\,A\leq\pi$, \deleted{et que $\mathrm{Ob}\,I=\bigcup_{r\in A}\mathrm{Ob}\,J_r$. (ou plus \unsure{généralt|généralement},} telles que $\bigcup_{r\in A}\mathrm{Ob}\,J_r$ soit cofinal dans $\mathrm{Ord}(I)$, i.e pour tt $i\in\mathrm{Ob}\,I$ \deleted{\unsure{et}} $\exists\,j\in\bigcup_{r\in A}\mathrm{Ob}\,J_r$, et une flèche $i\unsure{\text{---}|\longrightarrow}j$). Alors il existe \deleted{un} $r\in A$ tel\deleted{\unsure{le}} que $J_r$ soit cofinale dans $I$.

\deleted{\uline{\unsure{Proposition}}} \inserted{Théorème} 4.13.31. Soient $M,N$ des catégories accessibles, $f:M\unsure{\text{---}|\longrightarrow}N$ un foncteur accessible, \deleted{\unsure{$S$,}} $T\subset\mathrm{Ob}\,N$ une partie \deleted{accessible} de $N$, $\underline T\subset N$ la sous-catégorie str. pleine \unsure{correspondante} de $N$, \deleted{accessible dans $N$. Alors} $S=f^{-1}(T)$, $\underline S=f^{-1}(\underline T)$ la sous-catégorie pleine de $M$ associée. \unsure{Si} $T$ est accessible, alors $S$ est accessible\,; en d'autres termes, si $\underline T$ est accessible \unsure{en|dans} $N$, $\underline S$ est accessible dans $M$.

Cel\unsure{a} résulte de l'énoncé plus précis\unsure{:}

\uline{Corollaire} 4.13.32. Soient $\pi_0$ \inserted{et $\pi\unsure{>|\geq}\pi_0$} deux cardinaux ayant les propriétés suivantes :

\deleted{\unsure{1°} Il existe un \deleted{cardinal} ordinal limite $\alpha_0$, avec $\mathrm{Card}\,\alpha_0\leq\pi$, et tel que $\underline T$ soit $\alpha_0$-$\pi$-accessible dans $N$.}
% archive sheet number: 437
